\chapter{Słownik pojęć}
\label{cha:slownik}

Rozdział zawiera wyjaśnienie pojęć technicznych i branżowych pojawiających się w pracy. Zebrano je w jednym miejscu, aby ułatwić lekturę osobom nieznającym specyfiki branży kamieniarskiej lub terminologii związanej z tworzeniem aplikacji internetowych.

\begin{description}

\item[API] Interfejs programistyczny aplikacji; zbiór adresów, pod którymi jeden program udostępnia dane i operacje innemu programowi.

\item[Atrapa (mock)] Uproszczony zamiennik prawdziwego modułu, używany w testach. Zachowuje się tak, jak nakaże mu test, dzięki czemu można odtworzyć sytuacje trudne do wywołania na działającym systemie.

\item[Cross-site request forgery (CSRF)] Atak polegający na tym, że obca witryna wysyła żądanie do aplikacji, w której ofiara jest zalogowana, licząc na to, że przeglądarka automatycznie dołączy ciasteczka sesji.

\item[Cross-site scripting (XSS)] Atak polegający na wstrzyknięciu obcego kodu wykonywanego w przeglądarce ofiary, na przykład w celu wykradzenia tokenu sesji.

\item[Endpoint] Pojedynczy adres udostępniany przez serwer, pod którym można wykonać określoną operację.

\item[HttpOnly] Atrybut ciasteczka sprawiający, że jego wartość jest niedostępna dla kodu wykonywanego w przeglądarce, a przeglądarka dołącza je wyłącznie do żądań sieciowych.

\item[Inskrypcja] Napis wykonywany na tablicy nagrobnej, zawierający najczęściej imię, nazwisko oraz daty.

\item[JWT (JSON Web Token)] Podpisany token przenoszący informację o tożsamości użytkownika, weryfikowalny bez odpytywania bazy danych.

\item[jsdom] Programowa imitacja przeglądarki działająca w środowisku Node.js, pozwalająca uruchamiać testy interfejsu bez otwierania prawdziwego okna.

\item[Karta zamówienia] Projekt przygotowany przez klienta w konfiguratorze i zapisany w systemie. Nie jest jeszcze zamówieniem produkcyjnym; staje się nim po decyzji pracownika biura.

\item[Middleware (pośrednik)] Fragment kodu serwera wykonywany przed właściwą obsługą żądania, na przykład sprawdzający tożsamość użytkownika lub liczbę żądań z danego adresu.

\item[Migracja] Skrypt zmieniający strukturę bazy danych, zapisany w repozytorium, dzięki czemu identyczną bazę można odtworzyć na innym komputerze.

\item[Monter] Pracownik zakładu wykonujący montaż pomnika na cmentarzu. W systemie jest to jedna z trzech ról użytkownika.

\item[Ograniczanie liczby żądań (rate limiting)] Mechanizm odrzucający żądania, gdy z jednego adresu przyjdzie ich w krótkim czasie więcej, niż przewiduje limit.

\item[ORM] Biblioteka odwzorowująca tabele bazy danych na obiekty języka programowania. W projekcie świadomie z niej zrezygnowano.

\item[PostgREST] Usługa udostępniająca tabele bazy PostgreSQL bezpośrednio przez interfejs HTTP, wykorzystywana przez platformę Supabase.

\item[React] Biblioteka języka JavaScript służąca do budowania interfejsów użytkownika z niezależnych komponentów.

\item[REST] Styl budowy interfejsów sieciowych oparty na adresach zasobów i standardowych metodach protokołu HTTP.

\item[Row Level Security (RLS)] Mechanizm bazy PostgreSQL pozwalający określić, które wiersze tabeli widzi i może zmieniać dany użytkownik. Reguły działają w samej bazie, niezależnie od kodu aplikacji.

\item[Stela] Kształt tablicy nagrobnej w formie pionowej płyty o prostych krawędziach, wyższej niż szersza.

\item[Supabase] Platforma udostępniająca bazę PostgreSQL wraz z gotowym modułem kont i sesji oraz interfejsem HTTP do danych.

\item[Three.js] Biblioteka języka JavaScript służąca do wyświetlania grafiki trójwymiarowej w przeglądarce.

\item[TypeScript] Język programowania rozszerzający JavaScript o statyczne typowanie, wykrywające część błędów przed uruchomieniem programu.

\item[Wyrocznia testu] Wartość oczekiwana, z którą test porównuje wynik. Wyliczana niezależnie od sprawdzanego kodu, na przykład ze wzoru lub z treści wymagania.

\item[Wyzwalacz (trigger)] Procedura uruchamiana automatycznie przez bazę danych w reakcji na zapis, na przykład tworząca profil użytkownika przy zakładaniu konta.

\end{description}
